\documentclass[12pt,a4paper]{article}

% ==================== 基础宏包 ====================
\usepackage[UTF8]{ctex}              % 中文支持，UTF-8编码

% ==================== 页面设置 ====================
\usepackage{geometry}
\geometry{
    left=2.5cm,
    right=2.5cm,
    top=2.5cm,
    bottom=2.5cm
}

% ==================== 数学相关 ====================
\usepackage{amsmath}                 % 数学公式
\usepackage{amssymb}                 % 数学符号
\usepackage{amsthm}                  % 定理环境
\usepackage{mathtools}               % 数学工具增强

% ==================== 图表相关 ====================
\usepackage{graphicx}                % 图片支持
\usepackage{float}                   % 浮动体控制
\usepackage{caption}                 % 图表标题
\usepackage{subcaption}              % 子图支持
\usepackage{booktabs}                % 三线表
\usepackage{multirow}                % 表格多行合并
\usepackage{array}                   % 表格增强

% ==================== 其他常用宏包 ====================
\usepackage{hyperref}                % 超链接
\hypersetup{
    colorlinks=true,
    linkcolor=blue,
    filecolor=magenta,
    urlcolor=cyan,
    citecolor=green,
    bookmarksopen=true,
    bookmarksnumbered=true
}
\usepackage{cite}                    % 引用格式
\usepackage{enumerate}               % 列表环境
\usepackage{enumitem}                % 列表增强
\usepackage{listings}                % 代码高亮
\usepackage{xcolor}                  % 颜色支持
\usepackage{algorithm}               % 算法环境
\usepackage{algorithmicx}            % 算法排版
\usepackage{algpseudocode}           % 伪代码

% ==================== 代码高亮配置 ====================
\lstset{
    basicstyle=\ttfamily\small,
    keywordstyle=\color{blue},
    commentstyle=\color{gray},
    stringstyle=\color{red},
    numbers=left,
    numberstyle=\tiny\color{gray},
    stepnumber=1,
    numbersep=8pt,
    frame=single,
    breaklines=true,
    breakatwhitespace=false,
    showstringspaces=false,
    tabsize=4,
    captionpos=b
}

% ==================== 定理环境 ====================
\newtheorem{theorem}{定理}[section]
\newtheorem{lemma}[theorem]{引理}
\newtheorem{proposition}[theorem]{命题}
\newtheorem{corollary}[theorem]{推论}
\newtheorem{definition}{定义}[section]
\newtheorem{example}{例}[section]
\newtheorem{remark}{注}[section]

% ==================== 自定义命令 ====================
% 向量加粗
\renewcommand{\vec}[1]{\boldsymbol{#1}}
% 常用数学符号
\newcommand{\R}{\mathbb{R}}          % 实数集
\newcommand{\N}{\mathbb{N}}          % 自然数集
\newcommand{\Z}{\mathbb{Z}}          % 整数集
\newcommand{\C}{\mathbb{C}}          % 复数集
\newcommand{\Q}{\mathbb{Q}}          % 有理数集

% ==================== 标题信息 ====================
\title{\textbf{中文学术论文标准模板}}
\author{
    作者姓名\thanks{通讯作者邮箱: author@example.com} \\
    单位名称 \\
    地址信息
}
\date{\today}

% ==================== 正文开始 ====================
\begin{document}

\maketitle

% ==================== 摘要 ====================
\begin{abstract}
这是论文的中文摘要。摘要应简明扼要地概括论文的主要内容、方法、结果和结论。摘要应具有独立性和自含性，即不阅读全文就能获得论文的主要信息。摘要一般为200-300字，不分段。本模板提供了中文学术论文的标准结构和常用元素，包括标题、作者、摘要、关键词、正文各节、参考文献等。

\vspace{1em}
\noindent \textbf{关键词：} LaTeX；中文论文；学术写作；模板
\end{abstract}

% ==================== 目录（可选）====================
% \tableofcontents
% \newpage

% ==================== 正文部分 ====================

\section{引言}

引言部分应说明研究背景、研究意义、国内外研究现状、存在的问题以及本文的主要工作和创新点。

\subsection{研究背景}

研究背景介绍本研究领域的基本情况和发展历程。

\subsection{研究意义}

阐述本研究的理论意义和实际应用价值。

\subsection{文献综述}

回顾相关领域的研究现状。例如，引用文献\cite{latex}表明 LaTeX 是一个强大的排版系统。

\subsection{本文的主要工作}

简要说明本文的研究内容和创新点：
\begin{enumerate}
    \item 第一个创新点
    \item 第二个创新点
    \item 第三个创新点
\end{enumerate}

\section{相关工作}

详细介绍与本研究相关的已有工作，分析它们的优缺点。

\section{问题定义与数学模型}

\subsection{问题描述}

形式化描述要解决的问题。

\begin{definition}[问题定义]
给定输入空间 $\mathcal{X} \subseteq \R^n$ 和输出空间 $\mathcal{Y} \subseteq \R^m$，目标是找到映射函数 $f: \mathcal{X} \to \mathcal{Y}$。
\end{definition}

\subsection{数学模型}

建立问题的数学模型。例如，优化问题可表示为：

\begin{equation}
\begin{aligned}
    \min_{x \in \mathcal{X}} \quad & f(x) \\
    \text{s.t.} \quad & g_i(x) \leq 0, \quad i = 1, 2, \ldots, m \\
                      & h_j(x) = 0, \quad j = 1, 2, \ldots, p
\end{aligned}
\label{eq:optimization}
\end{equation}

其中 $f(x)$ 是目标函数，$g_i(x)$ 和 $h_j(x)$ 分别是不等式约束和等式约束。

\subsection{理论分析}

\begin{theorem}[收敛性定理]
在假设条件 A1-A3 下，算法~\ref{alg:example}~收敛到问题~\eqref{eq:optimization}~的最优解。
\end{theorem}

\begin{proof}
根据假设条件 A1-A3，我们可以构造 Lyapunov 函数并证明其单调递减性。具体证明过程涉及梯度下降的收敛性分析，由于篇幅限制，完整证明请参见附录~A。
\end{proof}

\section{方法}

\subsection{算法设计}

详细描述提出的算法或方法。

\begin{algorithm}[htbp]
\caption{示例算法}
\label{alg:example}
\begin{algorithmic}[1]
\Require 输入数据 $\mathcal{D} = \{(x_i, y_i)\}_{i=1}^n$
\Ensure 模型参数 $\theta$
\State 初始化参数 $\theta_0$
\For{$t = 1$ to $T$}
    \State 计算梯度 $g_t = \nabla_\theta \mathcal{L}(\theta_{t-1})$
    \State 更新参数 $\theta_t = \theta_{t-1} - \eta g_t$
    \If{收敛条件满足}
        \State \textbf{break}
    \EndIf
\EndFor
\State \Return $\theta_T$
\end{algorithmic}
\end{algorithm}

\subsection{实现细节}

描述算法的具体实现细节，包括参数设置、优化技巧等。

\section{实验}

\subsection{实验设置}

\subsubsection{数据集}

描述使用的数据集。表~\ref{tab:datasets}~总结了实验中使用的数据集统计信息。

\begin{table}[htbp]
    \centering
    \caption{数据集统计信息}
    \label{tab:datasets}
    \begin{tabular}{lcccc}
        \toprule
        \textbf{数据集} & \textbf{样本数} & \textbf{特征维度} & \textbf{类别数} & \textbf{来源} \\
        \midrule
        Dataset-A & 1,000 & 10 & 2 & [来源1] \\
        Dataset-B & 5,000 & 20 & 5 & [来源2] \\
        Dataset-C & 10,000 & 50 & 10 & [来源3] \\
        \bottomrule
    \end{tabular}
\end{table}

\subsubsection{评价指标}

使用以下评价指标：

\begin{itemize}
    \item \textbf{准确率 (Accuracy)}：$\text{Acc} = \frac{TP + TN}{TP + TN + FP + FN}$
    \item \textbf{精确率 (Precision)}：$\text{Prec} = \frac{TP}{TP + FP}$
    \item \textbf{召回率 (Recall)}：$\text{Rec} = \frac{TP}{TP + FN}$
    \item \textbf{F1分数}：$\text{F1} = \frac{2 \cdot \text{Prec} \cdot \text{Rec}}{\text{Prec} + \text{Rec}}$
\end{itemize}

\subsubsection{对比方法}

与以下基线方法进行对比：
\begin{enumerate}
    \item 方法A：简要描述
    \item 方法B：简要描述
    \item 方法C：简要描述
\end{enumerate}

\subsection{实验结果}

\subsubsection{主要结果}

表~\ref{tab:main-results}~展示了在不同数据集上的实验结果。

\begin{table}[htbp]
    \centering
    \caption{主要实验结果对比（\%）}
    \label{tab:main-results}
    \begin{tabular}{lcccc}
        \toprule
        \textbf{方法} & \textbf{Dataset-A} & \textbf{Dataset-B} & \textbf{Dataset-C} & \textbf{平均} \\
        \midrule
        方法A & 85.2 & 78.3 & 82.1 & 81.9 \\
        方法B & 87.4 & 80.5 & 84.3 & 84.1 \\
        方法C & 88.1 & 81.2 & 85.0 & 84.8 \\
        \textbf{本文方法} & \textbf{91.3} & \textbf{85.7} & \textbf{88.9} & \textbf{88.6} \\
        \bottomrule
    \end{tabular}
\end{table}

从表~\ref{tab:main-results}~可以看出，本文提出的方法在所有数据集上都取得了最好的性能。

\subsubsection{可视化结果}

图~\ref{fig:results}~展示了方法的可视化结果。

\begin{figure}[htbp]
    \centering
    % 取消注释以插入实际图片
    % \includegraphics[width=0.8\textwidth]{images/results.pdf}
    
    % 占位符
    \fbox{\parbox{0.8\textwidth}{\centering \vspace{4cm} 
    在此插入实验结果图 \\
    \texttt{images/results.pdf} 
    \vspace{4cm}}}
    
    \caption{实验结果可视化}
    \label{fig:results}
\end{figure}

\subsection{消融实验}

为了验证各个组件的有效性，进行了消融实验，结果如表~\ref{tab:ablation}~所示。

\begin{table}[htbp]
    \centering
    \caption{消融实验结果}
    \label{tab:ablation}
    \begin{tabular}{lcc}
        \toprule
        \textbf{模型变体} & \textbf{准确率 (\%)} & \textbf{F1分数 (\%)} \\
        \midrule
        完整模型 & 91.3 & 89.7 \\
        - 移除组件A & 88.5 & 86.2 \\
        - 移除组件B & 87.9 & 85.8 \\
        - 移除组件C & 89.2 & 87.4 \\
        \bottomrule
    \end{tabular}
\end{table}

\subsection{参数敏感性分析}

分析关键参数对模型性能的影响。

\section{讨论}

\subsection{结果分析}

深入分析实验结果，解释为什么提出的方法能取得更好的性能。

\subsection{局限性}

诚实地讨论方法的局限性和不足之处。

\subsection{未来工作}

展望后续研究方向：
\begin{itemize}
    \item 方向一
    \item 方向二
    \item 方向三
\end{itemize}

\section{结论}

总结全文的主要工作和贡献。本文提出了...，通过大量实验验证了...，结果表明...。

\section*{致谢}

感谢相关人员和机构的支持。本研究得到了...基金（项目编号：XXX）的资助。

% ==================== 参考文献 ====================
\begin{thebibliography}{99}

\bibitem{latex}
Lamport, Leslie. \textit{LaTeX: A Document Preparation System}. Addison-Wesley Professional, 1994.

\bibitem{ctex}
CTeX 开发组. \textit{CTeX 宏集手册}. 2021. \url{https://github.com/CTeX-org/ctex-kit}

\bibitem{example1}
作者1, 作者2. 论文标题[J]. 期刊名称, 年份, 卷号(期号): 页码.

\bibitem{example2}
Author A, Author B, Author C. Title of the paper[C]//Proceedings of Conference Name. Location: Publisher, Year: Pages.

\bibitem{example3}
作者. 书名[M]. 版次. 出版地: 出版社, 年份: 页码.

\end{thebibliography}

% ==================== 附录（可选）====================
\appendix

\section{附录A：推导细节}

提供正文中省略的详细推导过程。

\section{附录B：补充实验}

提供额外的实验结果和数据。

\end{document}
